\Article{45}
In applying the second method, it will be useful to distinguish two cases.

In the first case, where $a$ and $b$ have no common divisor, put $a\alpha+\beta b=1$, and let $k$ be the least positive residue of $\beta a-\alpha b$ modulo $p$. The identities
\[
\begin{aligned}
a(\beta a-\alpha b)&=\beta p-b(\alpha a+\beta b),\\
b(\beta a-\alpha b)&=-\alpha p+a(\alpha a+\beta b)
\end{aligned}
\]
then show that $ak\equiv-b,bk\equiv a(\mathrm{mod.}\ p)$. Thus, putting as above $ax+by=\xi$,

\SourceLeaf{126}%
$ay-bx=\eta$, we have $\eta\equiv k\xi$, $\xi\equiv-k\eta(\mathrm{mod.}\ p)$. All the numbers $\xi+\eta i$ to which simply least residues $x+yi$ correspond are therefore obtained by taking either for $\xi$, successively, the values $0,1,2,3\dotsfour p-1$, and for $\eta$ the least positive residues of the products $k\xi$ modulo $p$, or, in another order, for $\eta$ those values and for $\xi$ the least residues of the products $-k\eta$. From each $\xi+\eta i$ the corresponding $x+yi$ is then found by the formula
\[
x+yi=\frac{\xi+\eta i}{a-bi}=\frac{a\xi-b\eta}{p}+\frac{a\eta+b\xi}{p}\pdot i
\]

It is also clear that $\eta$, as $\xi$ increases by one, either increases by $k$ or decreases by $p-k$, and consequently $x+yi$ undergoes
\[
\begin{aligned}
&\text{either the change }\frac{a-kb}{p}+\frac{ak+b}{p}\pdot i\\
&\text{or this one }\frac{a-kb}{p}+b+\left(\frac{ak+b}{p}-a\right)i
\end{aligned}
\]
This observation serves to make the construction easier.

Finally, if absolutely least residues $x+yi$ are required, these instructions evidently change only in that $\xi$ must successively be given values between the bounds $-\frac12p$ and $+\frac12p$, while for $\eta$ we must take the absolutely least residues of the products $k\xi$. Here is a table of least residues for the modulus $5+2i$, arranged in this way:

\Needspace{15\baselineskip}
\begin{center}Simply least residues.\end{center}
{\small\setlength{\arraycolsep}{3pt}
\[
\begin{array}{r|r||r|r||r|r}
\xi+\eta i&x+yi&\xi+\eta i&x+yi&\xi+\eta i&x+yi\\\hline
0&0&10+25i&+5i&20+21i&+2+5i\\
1+17i&-1+3i&11+13i&+1+3i&21+9i&+3+3i\\
2+5i&+i&12+i&+2+i&22+26i&+2+6i\\
3+22i&+1+4i&13+18i&+1+4i&23+14i&+3+4i\\
4+10i&+2i&14+6i&+2+2i&24+2i&+4+2i\\
5+27i&-1+5i&15+23i&+1+5i&25+19i&+3+5i\\
6+15i&+3i&16+11i&+2+3i&26+7i&+4+3i\\
7+3i&+1+i&17+28i&+1+6i&27+24i&+3+6i\\
8+20i&+4i&18+16i&+2+4i&28+12i&+4+4i\\
9+8i&+1+2i&19+4i&+3+2i&&
\end{array}
\]
}

\SourceLeaf{127}%
\begin{center}Absolutely least residues.\end{center}
{\small\setlength{\arraycolsep}{3pt}
\[
\begin{array}{r|r||r|r||r|r}
\xi+\eta i&x+yi&\xi+\eta i&x+yi&\xi+\eta i&x+yi\\\hline
-14-6i&-2-2i&-4-10i&-2i&+5-2i&+1\\
-13+11i&-3+i&-3+7i&-1+i&+6-14i&+2-2i\\
-12-i&-2-i&-2-5i&-i&+7+3i&+1+i\\
-11-13i&-1-3i&-1+12i&-1+2i&+8-9i&+2-i\\
-10+4i&-2&0&0&+9+8i&+1+2i\\
-9-8i&-1-2i&+1-12i&+1-2i&+10-4i&+2\\
-8+9i&-2+i&+2+5i&+i&+11+13i&+1+3i\\
-7-3i&-1-i&+3-7i&+1-i&+12+i&+2+i\\
-6+14i&-2+2i&+4+10i&+2i&+13-11i&+3-i\\
-5+2i&-1&&&+14+6i&+2+2i
\end{array}
\]
}

The second case, where $a,b$ are not relatively prime, can easily be reduced to the preceding one. Let $\lambda$ be the greatest common divisor of $a,b$, and let $a=\lambda a'$, $b=\lambda b'$. Let $F$ denote any least residue modulo $\lambda$, with the modulus regarded as a complex number: that is, let it represent any number $x+yi$ such that $x,y$ lie either between the bounds $0$ and $\lambda$ or between $-\frac12\lambda$ and $+\frac12\lambda$, according as simply least or absolutely least residues are considered. Further, let $F'$ denote any least residue modulo $a'+b'i$. Then $(a'+b'i)F+F'$ will represent any least residue modulo $a+bi$, and a complete system of these residues results when all values of $F$ are combined with all values of $F'$.

\Article{46}
Two complex numbers are called relatively prime if they have no common divisors other than the units. Whenever such common divisors do exist, those whose norm is greatest are called greatest common divisors.

If the factorization of two given numbers into prime factors is available, the greatest common divisor is determined in exactly the same way as for real numbers (\emph{Disquiss. Ar.} Article 18). This also shows that every common divisor of two given numbers must divide their greatest common divisor found in this way. Since it is already evident that its three associates are also common divisors, there will always be four numbers, and no
\SourceLeaf{128}%
more, that are to be called greatest common divisors; their norm will be a multiple of the norm of any other common divisor.

If the factorization of two given numbers into simple factors is not available, the greatest common divisor is found by an algorithm similar to that for real numbers. Let $m,m'$ be the two given numbers, and form by repeated division the sequence $m'',m'''$, etc., so that $m''$ is the absolutely least residue of $m$ modulo $m'$, then $m'''$ the absolutely least residue of $m'$ modulo $m''$, and so on. Denoting the norms of $m,m',m'',m'''$, etc., respectively by $p,p',p'',p'''$, etc., $\frac{p''}{p'}$ will be the norm of the quotient $\frac{m''}{m'}$ and thus, by the definition of an absolutely least residue, certainly no greater than $\frac12$; the same holds for $\frac{p'''}{p''}$, etc. The positive real integers $p',p'',p'''$, etc., will therefore form a steadily decreasing sequence, so that a term $0$ must eventually be reached. Equivalently, in the sequence $m,m',m'',m'''$, etc., we shall eventually reach a term that divides the preceding term without a remainder. Let this be $m^{(n+1)}$, and put
\[
\begin{aligned}
m&=km'+m''\\
m'&=k'm''+m'''\\
m''&=k''m'''+m''''
\end{aligned}
\]
and so on, as far as
\[
m^{(n)}=k^{(n)}m^{(n+1)}
\]

Going through these equations in reverse order shows that $m^{(n+1)}$ divides each preceding term $m^{(n)}\dotsfour m'',m',m$; going through the same equations in forward order makes it clear that every common divisor of $m,m'$ also divides each subsequent term. The first conclusion shows that $m^{(n+1)}$ is a common divisor of $m,m'$; the second, that this divisor is greatest.

Whenever the last residue $m^{(n+1)}$ equals one of the four units $1,-1,i,-i$, this will show that $m$ and $m'$ are relatively prime.

\Needspace{22\baselineskip}
\Article{47}
If the equations of the preceding article, omitting the last, are combined so as to eliminate $m'',\allowbreak m''',\allowbreak m''''\dotsfour\allowbreak m^{(n)}$, an equation of the following form results:
\[
m^{(n+1)}=hm+h'm'
\]
\SourceLeaf{129}%
where $h,h'$ are integers; more specifically, if we use the notation introduced in \emph{Disquiss. Ar.} Article 27,
\[
\begin{aligned}
h&=\pm[k',k'',k'''\dotsfour k^{(n-1)}]\\
&=\pm[k^{(n-1)},k^{(n-2)}\dotsfour k'',k']\\
h'&=\mp[k,k',k'',k'''\dotsfour k^{(n-1)}]\\
&=\mp[k^{(n-1)},k^{(n-2)}\dotsfour k'',k',k]
\end{aligned}
\]
the upper or lower signs applying according as $n$ is even or odd. We state this theorem as follows:

\emph{A greatest common divisor of two complex numbers $m,m'$ can be put in the form $hm+h'm'$, with $h,h'$ integers.}

This evidently holds for the greatest common divisor to which the algorithm of the preceding article led, and also for its three associates. For these, the coefficients $h,h'$ must be replaced either by $hi,h'i$, by $-h,-h'$ or by $-hi,-h'i$.

Thus, whenever the numbers $m,m'$ are relatively prime, the following equation can be satisfied:
\[
1=hm+h'm'
\]

\Needspace{24\baselineskip}
For example, let the given numbers be $31+6i=m$, $11-20i=m'$. Here we find
\[
\begin{aligned}
k&=i,&m''&=+11-5i\\
k'&=+1-i,&m'''&=+5-4i\\
k''&=+2,&m''''&=+1+3i\\
k'''&=-1-2i,&m'''''&=+i\\
k''''&=+3-i
\end{aligned}
\]
and hence
\[
\begin{aligned}
\relax[k',k'',k''']&=-6-5i\\
[k,k',k'',k''']&=+4-10i
\end{aligned}
\]
and therefore
\[
m'''''=i=(6+5i)m+(4-10i)m'
\]
as well as
\[
1=(5-6i)m+(-10-4i)m'
\]
which is confirmed by carrying out the calculation.

\Article{48}
The foregoing has prepared everything required for the theory of congruences of the first degree in the arithmetic of complex numbers. But since this theory
\SourceLeaf{130}%
does not differ essentially from the one that applies to the arithmetic of real numbers and has been set out at length in the \emph{Disquisitiones Arithmeticae}, it will suffice to set down the main points here.

I. The congruence $mt\equiv1\ (\mathrm{mod.}\ m')$ is equivalent to the indeterminate equation $mt+m'u=1$. If this equation is satisfied by the values $t=h,u=h'$, the general solution of the congruence is expressed by $t\equiv h(\mathrm{mod.}\ m')$. The condition for solvability is that the modulus $m'$ and the coefficient $m$ have no common divisor.

II. The solution of the congruence $ax+b\equiv c(\mathrm{mod.}\ M)$, in the case where $a,M$ are relatively prime, depends on the solution of
\[
at\equiv1\ (\mathrm{mod.}\ M)
\]
If $t=h$ satisfies this congruence, the general solution of the former is contained in the formula
\[
x\equiv(c-b)h\ (\mathrm{mod.}\ M)
\]

\Needspace{12\baselineskip}\begingroup\postdisplaypenalty=10000\relax
III, The congruence $ax+b\equiv c(\mathrm{mod.}\ M)$, in the case where $a,M$ have a common divisor $\lambda$, is equivalent to
\[
\frac a\lambda\pdot x\equiv\frac{c-b}\lambda\ \left(\mathrm{mod.}\frac M\lambda\right)
\]
Thus, when $\lambda$ is taken to be the greatest common divisor of $a,M$, the solution of the given congruence is reduced to the preceding case. It is clear that the necessary and sufficient condition for solvability is that $\lambda$ also divide the difference $c-b$.\endgroup

\Article{49}
So far we have touched only on elementary matters, which nevertheless could not be omitted for the sake of the connection. In deeper investigations, the arithmetic of complex numbers resembles the arithmetic of real numbers in that more elegant and simpler theorems result when we admit only moduli that are prime numbers. Indeed, extending them to composite moduli usually involves more length than difficulty, and more work than ingenuity. The following discussion will therefore deal chiefly with prime moduli.

